\phantomsection
\section*{How to Keep This Book}
\addcontentsline{toc}{section}{How to Keep This Book}

This book is for a Catholic who keeps Advent, or means to, with the Roman books of 1962: the
Missal and the Breviary as the Code of Rubrics of 1960 orders them. It assumes faith and
goodwill and no knowledge of the older calendar.

It has three parts. The first, \emph{The Season}, says what the name Advent means, how the
Church came to keep the season, and how its four weeks are laid out. Read it in the days before
the First Sunday. The second, \emph{The Sundays and Days}, takes each Sunday, the Ember Days and
the last days before Christmas in turn. Read each section on the Saturday evening before its
Sunday and keep it by you through the week; each ends with counsel for that week. The third,
\emph{Keeping the Season}, gathers what the Fathers teach about the Lord's coming and our
watching, and turns it to practice: prayer, fasting and almsgiving, a rule of life, and
confession before Christmas. Read it once in the first week and come back to it as the weeks go
on.

The appendices are for daily use. Appendix~\ref{sec:daily} gives for each day of the season the
lesson from Isaiah that the Breviary reads, a short passage from the Fathers, and the
meditation St Alphonsus de Liguori assigns to that day, and from 17 December the O antiphon of
the evening; Alphonsus begins his Christmas novena of meditations on 16 December, and the table
follows him. Appendix~\ref{sec:prayers} prints the
season's prayers, each from a named edition, in Latin, in English, or in both.

The book is dated to Advent 2026. In the 1962 calendar the First Sunday of Advent is the Sunday
that falls on 30 November or nearest to it; in 2026 that is 29 November, and the season has
twenty-six days. The daily table is keyed to weeks and weekdays up to 16 December and to dates
from 17 December, so it serves a shorter or a longer Advent as well.

Three kinds of statement are kept apart. Where the book states law, it names the law, the
territory and the date, and it does not let the law of 1962 stand for the law of today. Where it
reports a custom, it says where and when the custom is found. Where it proposes a practice, it
calls the proposal its own counsel. Counsel is to be taken up with prudence, measured against
one's health, duties and state of life, and brought to a confessor or spiritual director where
one has one.
